% !TeX program = xelatex
% UTF-8 transcription of all nine handwritten source pages, in page order.
% Source wording and medical claims are retained without editorial correction.
% pic01.jpg = 20261012_102939897.jpg
% pic02.jpg = 20261012_102949454.jpg
% pic03.jpg = 20261012_102957295.jpg
% pic04.jpg = 20261012_103005603.jpg
% pic05.jpg = 20261012_103011957.jpg
% pic06.jpg = 20261012_103019516.jpg
% pic07.jpg = 20261012_103026950.jpg
% pic08.jpg = 20261012_103035077.jpg
% pic09.jpg = 20261012_103041999.jpg
% Compile directly with XeLaTeX, or input into a document using xepersian.
\ifdefined\settextfont
  \def\VRIStandalone{0}
\else
  \def\VRIStandalone{1}
  \documentclass[12pt,a4paper]{article}
  \usepackage[margin=2cm]{geometry}
  \usepackage{amsmath,amssymb}
  \usepackage{xepersian}
  \IfFontExistsTF{Amiri}{\settextfont{Amiri}}{\settextfont{FreeSerif}}
  \setlatintextfont{FreeSerif}
\fi
\providecommand{\en}[1]{\lr{#1}}
\newcommand{\vripage}[1]{\clearpage\section*{\en{VRI} ـ صفحهٔ #1}%
  {\footnotesize\en{No \ldots\quad Date \ldots}\par}}
\newcommand{\vrifooter}[1]{{\footnotesize\en{PAPCO}\quad #1\par}}
\ifnum\VRIStandalone=1 \begin{document}\fi

\vripage{۱}
\section*{\en{Common Viral Respiratory Infections -- 204}}
\begin{itemize}
\item انتقال بیماری می‌تواند قبل از بروز علائم بیماری رخ دهد.
\item ویروس‌های تنفسی کلاسیک: بیماری فقط در اپی‌تلیوم تنفسی ایجاد می‌شود.
  \begin{itemize}
  \item \en{Rhinoviruses}.
  \item در مقابل برخی ویروس‌ها درگیری شدید سیستمیک ایجاد می‌کنند.
  \item \en{SARS-CoV}، \en{Measles}.
  \end{itemize}
\item وقتی ویروس وارد جریان خون می‌شود علیه آن: \en{Life long immunity}.
\end{itemize}

\subsection*{انواع دسته‌بندی ویروس‌های تنفسی}
\begin{itemize}
\item بر اساس محل آناتومیکی درگیری در دستگاه تنفسی.
\item سندرم بالینی که ایجاد می‌کنند.
\item اتیولوژی.
\end{itemize}

\subsection*{محل آناتومیک درگیری در دستگاه تنفسی}
\begin{itemize}
\item تارهای صوتی مرز بین دستگاه تنفسی فوقانی و تحتانی را مشخص می‌کنند.
\item سلول‌های هدف در راه‌های تنفسی متفاوتند:
  \begin{itemize}
  \item سلول‌های اپی‌تلیال مژک‌دار یا بدون مژک.
  \item گابلت سل‌ها.
  \item سلول‌های \en{Clara}.
  \end{itemize}
\item پروسهٔ پاتولوژیک \en{wheezing} به دلیل انسداد یا باریک شدن راه‌های هوایی به دنبال انقباض عضلات صاف است.
\item در بافت ریه سلول مژک‌دار یا عضلهٔ صاف وجود ندارد بلکه پنوموسیت تیپ \en{I} و \en{II} وجود دارد.
  % The crossed-out word «غشایی» was replaced above the line by «پنوموسیت».
  \begin{itemize}
  \item پنومونی عفونت پنوموسیت‌ها است.
  \end{itemize}
\end{itemize}

\subsection*{سندرم‌های بالینی}
\begin{itemize}
\item مجاری هوایی فوقانی:
  \begin{itemize}
  \item رینیت (\en{Common Cold}).
  \item سینوزیت.
  \item اوتیت مدیا.
  \item کونژکتیویت.
  \item فارنژیت.
  \item تونسیلیت.
  \item لارینژیت.
  \end{itemize}
\item مجاری هوایی تحتانی:
  \begin{itemize}
  \item تراکئیت.
  \item برونشیت.
  \item برونشیولیت.
  \item پنومونی.
  \item آسم (اگزاسربیشن)، \en{Reactive airway} (\en{RAD}).
  \end{itemize}
\end{itemize}
\vrifooter{۱}

\vripage{۲}
\subsection*{برونشیولیت}
\begin{itemize}
\item طی آن هوا در ریه گیر می‌افتد، دشواری در بازدم وجود دارد: \en{wheezing}.
\item به دلیل التهاب یا عفونت برونشیول‌ها رخ می‌دهد.
  \begin{itemize}
  \item کوچک‌ترین راه‌های هوایی که بالاترین مقاومت را دارند.
  \end{itemize}
\end{itemize}

\subsection*{لارنگوتراکئیت}
\begin{itemize}
\item نام مرسوم: \en{Croup}؛ \en{Barky Cough}.
\item دشواری در دم.
\item التهاب در لارنگس، تراشه، برونش.
\end{itemize}

\subsection*{ویروس‌های تنفسی که می‌توانند علائم سیستمیک جدی ایجاد کنند}
\begin{itemize}
\item آنفولانزا.
\item سرخک.
\item کووید.
\item هانتاویروس (\en{HPS}).
\end{itemize}

\subsection*{ویروس‌های تنفسی در افراد ایمونوکامپتنت}
\begin{itemize}
\item کودکان بیشتر از بزرگسالان مبتلا به عفونت‌های ویروسی تنفسی می‌شوند.
\item شایع‌ترین علل عفونت جدی مجاری تنفسی تحتانی:
  \begin{itemize}
  \item آنفولانزا.
  \item \en{RSV}.
  \item متاپنوموویروس انسانی (\en{hMPV}).
  \end{itemize}
\item پاراآنفولانزا، آدنوویروس: آن‌ها نیز بیماری مزمن ایجاد می‌کنند.
\item رینوویروس‌ها (شایع‌ترین علت \en{CC}) می‌توانند باعث \en{LRTI} شوند.
\end{itemize}

\subsection*{ویروس‌های تنفسی در افراد ایمونوکامپرومایزد}
\begin{itemize}
\item ویروس‌های \en{latent phase}، \en{reactivating} دارند.
  \begin{itemize}
  \item \en{HSV}، \en{CMV}، آدنوویروس.
  \end{itemize}
\item در افراد با سرکوب ایمنی شایع‌ترین علت پنومونی \en{CMV} است.
\end{itemize}

\subsection*{آنفولانزا (ارتومیکسوویروس‌ها)}
\begin{itemize}
\item سندرم آنفولانزا: تب، میالژی، \en{fatigue}، گلودرد، سردرد، سرفه.
\item در ۱۰\% موارد یک عفونت باکتریال استافیلوکوکی / استرپتوکوکی روی آن سوار می‌شود.
\item ۴ سروتیپ از این ویروس وجود دارد.
  \begin{itemize}
  \item \en{A, C} گونه‌های متعددی را درگیر می‌کنند.
  \item \en{B} اختصاصاً انسان را درگیر می‌کند.
  \item \en{D} دام را درگیر می‌کند، در انسان بیماری‌زایی ندارد.
  \item تیپ \en{A} بیماری‌زاترین و مسئول شدیدترین بیماری را ایجاد می‌کند.
  \end{itemize}
\end{itemize}
\vrifooter{۲}

\vripage{۳}
\subsection*{\en{hRSV}}
\begin{itemize}
\item از خانوادهٔ پارامیکسوویروس‌ها می‌باشد.
\item اکثر افراد ظرف سال‌های ابتدایی زندگی مبتلا می‌شوند ولی ابتلای مجدد در بزرگسالی نیز شایع است.
\item یکی از انتقال‌پذیرترین ویروس‌ها در انسان است.
\item اپیدمی‌های سالانه در اکتبر، نوامبر، مارس، در مناطق استوایی.
\item یکی از شایع‌ترین علل عفونت \en{LRT} در افراد مسن و کودکان.
\item از جمله مهم‌ترین علل بستری در دو سر طیف سنی است.
\item عفونت مجدد با این ویروس شایع است.
\item عفونت مجدد معمولاً به صورت یک بیماری خفیف \en{URT} می‌باشد.
\item افراد با نقص ایمنی در معرض خطر عفونت مزمن قرار دارند.
  \begin{itemize}
  \item خصوصاً افراد \en{BMT}.
  \end{itemize}
\item عفونت اولیه در کودکان به طور شایع به صورت برونشیولیت (ویزینگ) تظاهر می‌شود.
\item ارتباط قوی بین ابتلا به \en{hRSV} در سنین کودکی و ابتلا به آسم وجود دارد.
\item در بزرگسالان عامل:
  \begin{itemize}
  \item \en{COPD Exacerbation}.
  \item \en{AECB}؛ اگزاسربیشن حاد برونشیت مزمن.
  \end{itemize}
\end{itemize}

\subsection*{\en{hMPV}}
\begin{itemize}
\item عفونت اولیه معمولاً در ابتدای کودکی رخ می‌دهد.
\item عفونت مجدد در طول زندگی شایع است.
\item مشابه \en{hRSV} است.
\item هم عفونت \en{URT} و هم \en{LRT} ایجاد می‌کند.
\item انتقال‌پذیری آن از \en{hRSV} کمتر است.
\item علائم درگیری \en{LRT} در آن مشابه سایر پارامیکسوویروس‌ها است.
  \begin{itemize}
  \item سرفه، کوریزا، ویزینگ.
  \end{itemize}
\item نقش مهمی در \en{Exac.} آسم، \en{COPD} دارد.
\end{itemize}

\subsection*{\en{hPIV}}
\begin{itemize}
\item ویروس پاراآنفولانزا انسانی: خانوادهٔ پارامیکسوویروس‌ها.
\item \en{hPIV3} (سویهٔ ویروس نمونه [خوانش نامطمئن]) شایع‌ترین عامل عفونت‌شده در این خانواده است.
  % The parenthetical handwritten wording is uncertain; see pic03.jpg.
\item عفونت‌های مجدد متعدد در طول زندگی می‌تواند رخ دهد.
  \begin{itemize}
  \item اگرچه که عفونت‌های ثانویه خفیف یا بدون علامت هستند.
  \end{itemize}
\item عفونت اولیه در کودکان به صورت لارنگوتراکئیت (\en{Croup}) رخ می‌دهد.
  \begin{itemize}
  \item عفونت‌های ثانویه معمولاً صرفاً \en{URT} را درگیر می‌کنند.
  \end{itemize}
\item تشخیص با \en{RT-PCR} یا کشت سلولی است.
\end{itemize}
\vrifooter{۳}

\vripage{۴}
\subsection*{سرخک (\en{Measles})}
\begin{itemize}
\item یکی از اعضای خانوادهٔ پارامیکسوویروس‌ها است.
\item عفونت سیستمیک به نام \en{Rubeola} ایجاد می‌کند.
  \begin{itemize}
  \item می‌تواند عفونت تنفسی صرف ایجاد کند.
  \end{itemize}
\item واگیرترین ویروس تنفسی در انسان است.
\item با واکسیناسیون قابل پیشگیری است.
\item اگر نرخ واکسیناسیون زیر ۹۰--۹۵\% برود بیماری بالینی دیده می‌شود.
\item علائم شامل:
  \begin{itemize}
  \item ۳ \en{C} در تب شدید: سرفه + کوریزا + کونژکتیویت.
  \item ویرمی ایجاد می‌کند.
  \item راش ماکولوپاپولر منتشر / \en{Koplik's spots}: در مخاط دهان.
  \item در وجود راش تیپیک و تب مشاهدهٔ آن تشخیص را قطعی می‌کند.
  \end{itemize}
\end{itemize}

\subsection*{پیکورناویروس‌ها}
\begin{itemize}
\item انتروویروس‌ها، رینوویروس‌ها.
\item رینوویروس‌ها \en{year-round}، انتروویروس‌ها در تابستان در مناطق گرم.
\end{itemize}

\subsection*{رینوویروس‌ها}
\begin{itemize}
\item شایع‌ترین ویروس‌های بیماری‌زا در انسان هستند.
\item شایع‌ترین عامل \en{Common Cold} هستند: ۵۰\% موارد.
\item عارضهٔ \en{Common Cold} در کودکان \en{Otitis}، در بزرگسالان سینوزیت است.
\item با تشدید بیماری \en{RAD} در کودکان همراهی دارد.
  \begin{itemize}
  \item و آسم در بزرگسالان.
  \end{itemize}
\item رینوویروس \en{C} عامل شدیدترین بیماری‌ها است: پنومونی، \en{COPD-Ex.}
\item تا مدت‌ها پس از بیماری \en{RT-PCR} مثبت می‌باشد.
\end{itemize}

\subsection*{انتروویروس‌ها}
\begin{itemize}
\item انتروویروس‌های \en{non-polio} شایع هستند ولی عفونت آن‌ها معمولاً بدون علامت است.
\item ولی اپیدمی‌های بیماری‌زا می‌توانند ایجاد کنند.
\item کوکساکی‌ویروس‌ها در این دسته جای دارند.
\item ضایعات دهانی ایجاد می‌کنند: بیماری دست ـ پا ـ دهان.
\item فارنژیت ناشی از آن‌ها معمولاً با هرپانژین (زخم در کام) همراه است.
  \begin{itemize}
  \item تب، گلودرد، دشواری در بلع.
  \item اپیدمی معمولاً در کودکان کم‌سن در تابستان رخ می‌دهد.
  \end{itemize}
\item اکوویروس‌ها:
  \begin{itemize}
  \item در کودکان با عفونت \en{URT} در فصل تابستان.
  \item لارنگوتراکئیت (\en{Croup}) ایجاد می‌کنند.
  \item اپیدمی پلورودینیا (\en{pleurodynia}): تب + درد عضلات قفسهٔ سینه.
  \end{itemize}
\end{itemize}
\vrifooter{۴}

\vripage{۵}
% Continues the enterovirus material from source page 4.
\begin{itemize}
\item انتروویروس \en{D68}.
\item ویزینگ + سندرم شبه پولیو: \en{Acute Flaccid Myelitis} در کودکان.
\end{itemize}

\subsection*{آدنوویروس‌ها}
\begin{itemize}
\item هم انسان و هم حیوانات را درگیر می‌کنند.
\item برخی سروتیپ‌ها می‌توانند بیماری مزمن بدون علامت ایجاد کنند.
\item گونهٔ \en{B, C} آن در انسان عفونت دستگاه تنفسی ایجاد می‌کنند.
\item می‌توانند \en{CC}، \en{Croup}، برونشیولیت، پنومونی ایجاد کنند.
\item کونژکتیویت در زمینهٔ عفونت گونه‌های \en{B} و \en{D} رخ می‌دهد.
\item سندرم بالینی شایع آدنوویروس‌ها: تب فارنگوکونژکتیوال.
\item گاستروانتریت هم می‌توانند ایجاد کنند: گونهٔ \en{F}.
\item افراد نقص ایمنی بسیار مستعد عفونت آدنوویروسی شدید هستند.
\item می‌تواند \en{ARDS} ایجاد کند: سربازان آمریکایی.
  \begin{itemize}
  \item واکسن خوراکی آدنوویروس برای ۱۷--۵۰ سال.
  \end{itemize}
\end{itemize}

\subsection*{کروناویروس‌ها}
\begin{itemize}
\item برخلاف سایر کروناها، ویروس \en{SARS} علائم تنفسی فوقانی ایجاد نمی‌کند.
  \begin{itemize}
  \item سرفه و دیس‌پنه در اکثر بیماران دیده می‌شود.
  \item معمولاً تب، میالژی، ملز، \en{chills}، \en{rigors}، اسهال آبکی.
  \end{itemize}
\end{itemize}

\subsection*{هرپس‌ویروس‌ها}
\begin{itemize}
\item فارنژیت هرپسی: استوماتیت اولسراتیو حاد، فارنژیت اولسراتیو.
\item ۹۰\% ضایعات دهانی هرپسی مربوط به \en{HSV-1} است.
\item بیماری اولیهٔ دهانی در کودکان می‌تواند شدید باشد.
\item درگیری \en{EBV}: فارنژیت اگزوداتیو حاد و تحت‌حاد.
  \begin{itemize}
  \item گاهی التهاب لوزه آنقدر شدید می‌شود که راه هوایی را تهدید می‌کند.
  \end{itemize}
\item ویروس‌های خانوادهٔ هرپس \en{CMV}، \en{VZV}، \en{EBV}، \en{HHV} در افراد نقص ایمنی خصوصاً \en{BMT} می‌توانند بیماری‌های بسیار خطرناک ایجاد کنند.
\end{itemize}

\subsection*{پاربوویروس‌ها}
\begin{itemize}
\item خصوصاً در کودکان کم‌سن بستری علت عفونت تنفسی ایجاد می‌شوند.
\item انتقال یافتن آن در افراد سالم بی‌علامت در \en{RT-PCR} نیز بالا است.
\end{itemize}

\subsection*{رتروویروس‌ها}
\begin{itemize}
\item \en{HIV}: عفونت اولیهٔ خود می‌تواند فارنژیت با زخم‌های دهانی و لنفادنوپاتی ایجاد کند.
\end{itemize}
\vrifooter{۵}

\vripage{۶}
\subsection*{اپیدمیولوژی}
\begin{itemize}
\item سن یکی از مهم‌ترین عوامل تعیین‌کننده برای عفونت‌های تنفسی علامت‌دار است.
\item در اغلب موارد عفونت اولیه از عفونت‌های ثانویه شدیدتر است.
\item اکثر موارد عفونت‌های اولیه در سال‌های ابتدایی زندگی رخ می‌دهند.
\item مواجهه با کودکان یک ریسک فاکتور برای عفونت‌های ثانویه است.
\item فصل:
  \begin{itemize}
  \item اغلب ویروس‌های تنفسی شایع در زمستان بیشتر باعث بیماری می‌شوند.
  \item کرونا از فصل پیروی نمی‌کند بلکه تغییرات \en{Ag} در خصوص آن مهم است.
  \end{itemize}
\item ریسک فاکتورها:
  \begin{itemize}
  \item بیماری زمینه‌ای ریوی بسیار مهم است: \en{COPD}.
  \item سایر بیماری‌های زمینه‌ای: قلبی عروقی.
  \item دخانیات، وضعیت اجتماعی اقتصادی ضعیف.
  \item جنس مرد.
  \item چاقی: وضعیت التهابی مزمن.
  \end{itemize}
\end{itemize}

\subsection*{اپروچ بالینی به بیمار با عفونت‌های شایع ویروسی تنفسی}
\begin{itemize}
\item اصول کلی مراقبت، مراقبت حمایتی است.
\item هایپوکسی با \en{$\mathrm{O}_2$ therapy} و نارسایی تنفسی با \en{MV} حمایت می‌شود.
\item مدیریت مایعات به دلیل ریسک بالای دهیدریشن.
\item باید تصمیم بگیریم که آیا نیاز به درمان ضد ویروسی / ضد باکتری وجود دارد یا نه.
\item درمان‌های ضد ویروسی فقط زمانی مؤثر هستند که در ابتدای بیماری مصرف شوند.
\end{itemize}

\subsection*{تظاهرات بالینی}
\begin{itemize}
\item \en{CC}: احتقان بینی، عطسه، رینوره، سرفه، گلودرد.
\item لارینژیت معمولاً با دیس‌فونی، \en{hoarseness} همراه است.
\item برونشیت حاد:
  \begin{itemize}
  \item سرفهٔ خشک یا پروداکتیو کمتر از ۳ هفته.
  \item معمولاً در زمستان.
  \item در فقدان علائم و نشانه‌های پنومونی در \en{CXR}.
  \end{itemize}
\item برونشیولیت: بیماری حاد + ویزینگ + شواهد درگیری \en{URT}.
  \begin{itemize}
  \item بیشتر در شیرخواران و کودکان در فصول سرد دیده می‌شود.
  \end{itemize}
\item پنومونی: سرفه + خلط + دیس‌پنه + درد قفسهٔ سینه.
  \begin{itemize}
  \item علائم سیستمیک: تب، خستگی، تعریق + سردرد + میالژی.
  \end{itemize}
\end{itemize}

استاندارد طلایی تشخیص عفونت ویروسی تنفسی: ایزوله کردن ویروس.

\en{CXR} برای هر بیمار مشکوک به پنومونی ضروری است.
\vrifooter{۶}

\vripage{۷}
\subsection*{واکسن}
\begin{itemize}
\item آنفولانزا سالانه نیاز به تجدید.
\item آدنوویروس سروتیپ ۴ و ۷: برای نیروهای نظامی.
\item \en{RSV} برای افراد مسن:
  \begin{itemize}
  \item تک‌دوز برای افراد بالای ۶۰ سال.
  \item تک‌دوز برای زنان باردار بین ۳۲--۳۶ هفتگی.
  \end{itemize}
\end{itemize}

\subsection*{\en{Palivizumab}}
\begin{itemize}
\item آنتی‌بادی مونوکلونال ضد \en{RSV}.
\item برای جلوگیری از بستری کودکان.
\end{itemize}

\subsection*{\en{Ribavirin}}
\begin{itemize}
\item جلوگیری از تکثیر \en{RNA} ویروس \en{RSV}، آدنوویروس، هانتاویروس، سرخک، \en{PIV}، آنفولانزا.
\end{itemize}

\subsection*{\en{Pleconaril}}
\begin{itemize}
\item خوراکی برای درمان عفونت رینوویروس.
\item کاهش ترشح مخاط و علائم بیماری.
\end{itemize}

\subsection*{\en{Aciclovir}}
\begin{itemize}
\item آنالوگ گوانین.
\item درمان ویروس‌های خانوادهٔ هرپس.
\end{itemize}

\subsection*{مهارکننده‌های نورآمینیداز}
\begin{itemize}
\item مؤثر بر آنفولانزا \en{A} و \en{B}.
\item اوسلتامیویر: پیشگیری و درمان آنفولانزا بدون عارضه.
  \begin{itemize}
  \item برخی آنفولانزاها به آن مقاوم هستند: زانامیویر می‌دهیم.
  \end{itemize}
\item پرامیویر: به فرم وریدی، آنفولانزا بدون عارضه.
\end{itemize}

\section*{\en{The Common Cold}}
\begin{itemize}
\item یک سندرم بالینی است: \en{Mild Upper-Respiratory Viral Illness}.
\item دفعات ابتلا سالیانه:
  \begin{itemize}
  \item ۵--۷ نوبت در کودکان پیش از مدرسه.
  \item ۲--۳ نوبت در بزرگسالان.
  \end{itemize}
\item ویروس‌ها: رینوویروس‌ها (۳۰--۵۰\% موارد)، کروناویروس‌ها (۱۰--۱۵\%)، \en{RSV}، پاراآنفولانزا، آدنوویروس، انتروویروس.
\end{itemize}

\subsection*{اپیدمیولوژی}
\begin{itemize}
\item رینوویروس و پاراآنفولانزا در پاییز و اواخر بهار طغیان می‌کنند.
\item \en{RSV} و کرونا معمولاً در زمستان و بهار طغیان می‌کنند.
\item انتروویروس بیشتر در تابستان.
\item آدنوویروس فصلی نیست.
\end{itemize}
\vrifooter{۷}

\vripage{۸}
\begin{itemize}
\item ویروس‌های عامل سرماخوردگی می‌توانند تا ۲ ساعت روی پوست دست زنده بمانند.
\item ویروس‌های عامل سرماخوردگی در بزاق وجود ندارند، از طریق بزاق منتقل نمی‌شوند.
\item اوج انتقال ویروس معمولاً در روز اول و دوم بیماری است ولی می‌تواند تا دو هفته ادامه داشته باشد.
\end{itemize}

\subsection*{ریسک فاکتورهای ابتلا}
\begin{itemize}
\item مواجهه با کودکان در مهدکودک.
\item استرس سایکولوژیک.
\item فعالیت فیزیکی متوسط شانس ابتلا را کاهش می‌دهد.
\item کمبود خواب و اختلال خواب ریسک ابتلا را افزایش می‌دهد.
\item هوای سرد تأثیری بر استعداد ابتلا به \en{CC} ندارد.
\end{itemize}

\subsection*{تظاهرات بالینی}
\begin{itemize}
\item رینیت و احتقان بینی شایع‌ترین علامت‌های بالینی هستند.
\item گلودرد، سرفه، ضعف و خستگی.
\item تب در بزرگسالان نادر ولی در کودکان شایع است.
\item کونژکتیویت می‌تواند در برخی رخ دهد.
\item سمع \en{rale} و شواهد درگیری \en{LRT} وجود ندارد.
\end{itemize}

\begin{itemize}
\item رنگ ترشحات بینی بیانگر عفونی بودن آن نیست.
\item دورهٔ کمون برای اکثر ویروس‌های \en{CC}، ۲۴--۷۲ ساعت است.
\item طول دورهٔ بیماری معمول ۳--۷ روز است، تا نهایتاً ۲ هفته در افراد سیگاری.
\item سرفه می‌تواند ۴ هفته‌ها بعد از رفع بیماری ادامه پیدا کند.
\item بررسی اتیولوژیک روتین ضروری ندارد.
\item شدینگ ویروس در فرد بدون علامت خصوصاً در رینوویروس‌ها شایع است و در ۶\% افراد وجود دارد.
\end{itemize}

\subsection*{\en{DDX}}
\begin{itemize}
\item رینیت آلرژیک / فصلی.
\item فارنژیت / تونسیلیت باکتریال.
\item رینوسینوزیت باکتریال حاد.
\item آنفولانزا.
\item پرتوسیس.
\end{itemize}
\vrifooter{۸}

\vripage{۹}
\begin{itemize}
\item افتراق از رینیت: گلودرد و سرفه.
\item افتراق از تونسیلیت باکتریال: وجود رینوره و احتقان بینی.
\item افتراق از پرتوسیس: سرفهٔ شدید بیش از دو هفته.
\item افتراق از کووید: از دست دادن حس بویایی (\en{Anosmia}) و چشایی (\en{Ageusia}) به نفع کووید است.
\end{itemize}

\subsection*{درمان‌های علامتی}
\begin{itemize}
\item ضد درد: استامینوفن و \en{NSAID} به یک اندازه مؤثر هستند.
\item آنتی‌هیستامین / دکونژستان: ترکیبشان از هر کدام جداگانه مؤثرتر است.
  \begin{itemize}
  \item آنتی‌هیستامین به تنهایی معمولاً مؤثر نیست.
  \item نسل \en{I} باید بدهیم.
  \end{itemize}
\item اسپری بینی ایپراتروپیوم: رینوره و عطسه را کنترل می‌کند ولی بر احتقان بی‌تأثیر است.
\item دکونژستان (سودوافدرین) به تنهایی تأثیر کمی دارند.
  \begin{itemize}
  \item در فرم موضعی حداکثر تا ۳ روز می‌دهیم.
  \end{itemize}
\item \en{vit C}، \en{Zinc}، اکیناسه: بی‌تأثیر هستند.
\item بخور سرد و گرم: بی‌تأثیر.
\end{itemize}
\vrifooter{۹}

\ifnum\VRIStandalone=1
  \let\VRIEnd\enddocument
\else
  \let\VRIEnd\relax
\fi
\VRIEnd
